\documentclass[11pt]{article}
\usepackage[margin=1in]{geometry}
\usepackage{amsmath,amssymb}
\usepackage{mathtools}
\usepackage{parskip}

\allowdisplaybreaks

\title{STAT GU4204 --- Homework 3}
\author{Jonah Aden}
\date{Due: 04/28/2026}

\newcommand{\E}{\mathbb{E}}
\newcommand{\Var}{\mathrm{Var}}
\newcommand{\Cov}{\mathrm{Cov}}
\newcommand{\iid}{\stackrel{\text{iid}}{\sim}}
\newcommand{\CRLB}{\mathrm{CRLB}}

\begin{document}
\maketitle

%----------------------------------------------------------------
\section*{8.5.2(a)}
Let $X_1,\dots,X_n \iid N(\mu,\sigma^2)$ with $\mu,\sigma^2$ unknown, and $(S_n')^2 = \frac{1}{n-1}\sum(X_i-\bar X_n)^2$.

By Fisher--Cochran, $\bar X_n \perp (S_n')^2$, so
\[
T = \sqrt{n}\,\frac{\bar X_n - \mu}{S_n'} \sim t_{n-1}
\]
is a pivot. For $c = T_{n-1,(1+\gamma)/2}$, symmetry gives $\Pr(-c \le T \le c)=\gamma$. Inverting:
\[
\boxed{\left(\bar X_n - T_{n-1,(1+\gamma)/2}\frac{S_n'}{\sqrt{n}},\;\; \bar X_n + T_{n-1,(1+\gamma)/2}\frac{S_n'}{\sqrt{n}}\right)}
\]
is a $\gamma$ confidence interval for $\mu$.

\bigskip
%----------------------------------------------------------------
\section*{8.5.4}
Let $X_1,\dots,X_n \iid N(\mu,\sigma^2)$ with both parameters unknown.

Since $(n-1)(S_n')^2/\sigma^2 \sim \chi^2_{n-1}$, this is a pivot. Using equal-tail cutoffs $a=\chi^2_{n-1,\alpha/2}$, $b=\chi^2_{n-1,1-\alpha/2}$ with $\alpha=1-\gamma$, we have $\Pr(a \le V \le b)=\gamma$. Inverting for $\sigma^2$:
\[
\boxed{\left(\frac{(n-1)(S_n')^2}{\chi^2_{n-1,1-\alpha/2}},\;\;\frac{(n-1)(S_n')^2}{\chi^2_{n-1,\alpha/2}}\right), \quad \alpha=1-\gamma}
\]

\bigskip
%----------------------------------------------------------------
\section*{8.7.11(a)}
Let $X_1,\dots,X_n \iid N(\mu,\sigma^2)$ and $T_c = c\sum_{i=1}^n(X_i-\bar X_n)^2$.

Since $\sum(X_i-\bar X_n)^2/\sigma^2 \sim \chi^2_{n-1}$, we have $\E\bigl[\sum(X_i-\bar X_n)^2\bigr]=(n-1)\sigma^2$, so
\[
\E[T_c] = c(n-1)\sigma^2.
\]
Setting $\E[T_c]=\sigma^2$ requires $c(n-1)=1$, so $\displaystyle\boxed{c = \frac{1}{n-1}}$.

\bigskip
%----------------------------------------------------------------
\section*{8.8.2}
Let $f(x\mid\theta)=\theta e^{-\theta x}$ for $x>0$.

The log-likelihood is $\ell = \log\theta - \theta x$, giving $\dot\ell = 1/\theta - x$ and $\ddot\ell = -1/\theta^2$. Therefore
\[
I(\theta) = -\E[\ddot\ell] = \boxed{\dfrac{1}{\theta^2}}.
\]

\bigskip
%----------------------------------------------------------------
\section*{8.8.7}
Let $X_1,\dots,X_n \iid N(\mu,\sigma^2)$ with $\mu$ known, and $\hat\sigma^2 = \frac{1}{n}\sum(X_i-\mu)^2$.

\textit{Unbiasedness.} Since $\sum(X_i-\mu)^2/\sigma^2 \sim \chi^2_n$, we get $\E[\hat\sigma^2]=\sigma^2$.

\textit{Variance.} Using $\Var(\chi^2_n)=2n$:
\[
\Var(\hat\sigma^2) = \frac{\sigma^4}{n^2}\cdot 2n = \frac{2\sigma^4}{n}.
\]

\textit{Fisher information for $\theta=\sigma^2$.} From $\ell = -\frac{1}{2}\log(2\pi\theta)-\frac{(x-\mu)^2}{2\theta}$:
\[
\ddot\ell = \frac{1}{2\theta^2} - \frac{(x-\mu)^2}{\theta^3} \implies I(\sigma^2) = \frac{1}{2\sigma^4}.
\]

\textit{Efficiency.} $\displaystyle\CRLB = \frac{1}{nI(\sigma^2)} = \frac{2\sigma^4}{n} = \Var(\hat\sigma^2)$.

Since $\hat\sigma^2$ is unbiased and attains the CRLB, it is efficient. $\square$

\bigskip
%----------------------------------------------------------------
\section*{8.8.17}
Let $T$ be unbiased for $g(\theta)$ and $U_n=\sum_{i=1}^n\dot\ell(X_i,\theta)$. Under regularity conditions, $\E[U_n]=0$, $\Var(U_n)=nI(\theta)$, and differentiating $\E[T]=g(\theta)$ under the integral gives $\Cov(T,U_n)=g'(\theta)$.

Applying Cauchy--Schwarz to the mean-zero variables $T-g(\theta)$ and $U_n$:
\[
[g'(\theta)]^2 = [\Cov(T,U_n)]^2 \le \Var(T)\cdot nI(\theta),
\]
which is the CRLB. Equality holds iff $T-g(\theta) = u(\theta)U_n$ a.s. Taking covariance with $U_n$:
\[
g'(\theta) = u(\theta)\cdot nI(\theta) \implies u(\theta) = \frac{g'(\theta)}{nI(\theta)},
\]
and then $\Var(T)=u(\theta)^2\cdot nI(\theta)=[g'(\theta)]^2/(nI(\theta))=\CRLB$. $\square$

\bigskip
%----------------------------------------------------------------
\section*{9.1.1(a)}
Let $X_1,\dots,X_n \iid N(\mu,1)$; reject $H_0:\mu=0$ when $\bar X_n > c$.

Under $H_0$, $\bar X_n \sim N(0,1/n)$:
\[
\alpha(c) = \Pr_{\mu=0}(\bar X_n > c) = 1-\Phi(\sqrt{n}\,c).
\]
Under $H_1:\mu=1$, $\bar X_n \sim N(1,1/n)$:
\[
\beta(c) = \Pr_{\mu=1}(\bar X_n \le c) = \Phi\!\bigl(\sqrt{n}(c-1)\bigr).
\]
\[
\boxed{\alpha(c)=1-\Phi(\sqrt{n}\,c),\qquad \beta(c)=\Phi\!\bigl(\sqrt{n}(c-1)\bigr).}
\]

\end{document}
